% Source PDF: source/普通高考/2009/2009安徽文.pdf, page 2, question 12.
% Program-flow topology: 开始→a=1→a=2a+1→a>100?;
% 是→输出a→结束；否→左侧回边→a=2a+1入口主干。
% The source bitmap img/2009/anhui_liberal/q12_fig1.png was redrawn as vector TikZ.

\begin{tikzpicture}[
  x=1cm,
  y=0.9cm,
  >=Stealth,
  line width=0.45pt,
  line cap=round,
  line join=round,
  flowtext/.style={font=\normalsize\setlength{\parindent}{0pt}\noindent,
    anchor=center,align=center},
  every node/.style={flowtext},
  flowbox/.style={flowtext,draw,minimum height=0.68cm,inner xsep=4pt,inner ysep=2.5pt},
  branchlabel/.style={flowtext,inner sep=0pt},
  terminal/.style={flowbox,rounded corners=0.20cm,minimum width=1.35cm},
  io/.style={flowbox,shape=trapezium,trapezium left angle=72,
    trapezium right angle=108,minimum width=2.40cm,minimum height=0.76cm},
  decision/.style={flowbox,shape=diamond,aspect=3.2,minimum width=3.85cm,
    minimum height=1.00cm,inner xsep=4pt,inner ysep=2pt}
]
  \node[terminal] (start) at (0,8.70) {开始};
  \node[flowbox,minimum width=2.05cm] (init) at (0,7.30) {$a=1$};
  \node[flowbox,minimum width=3.30cm] (update) at (0,5.70) {$a=2a+1$};
  \node[decision] (test) at (0,4.10) {$a>100?$};
  \node[io] (output) at (0,2.55) {输出 $a$};
  \node[terminal] (finish) at (0,1.10) {结束};

  \draw[->] (start.south) -- (init.north);
  \draw[->] (init.south) -- (update.north);
  \draw[->] (update.south) -- (test.north);
  \draw[->] (test.south) -- node[branchlabel,right=2pt] {是} (output.north);
  \draw[->] (output.south) -- (finish.north);

  % 否 branch: independent return arrow enters the main stem before a=2a+1.
  \coordinate (loopJoin) at (0,6.55);
  \draw (test.west) -- node[branchlabel,above=2pt] {否}
    (-2.60,4.10) -- (-2.60,6.55);
  \draw[->] (-2.60,6.55) -- (loopJoin);
\end{tikzpicture}
